\section{Poruszanie się po labiryncie}
\subsection{Ogólna zasada}
Gra \textit{Labirynt śmierci} nie posiada planszy w ścisłym tego słowa znaczeniu, gracze tworzą ją w trakcie gry, konstruując stopniowo labirynt poprzez układanie pasujących do siebie żetonów, oznaczających korytarze i komnaty.
Żeton reprezentujący oddział musi zawsze być umieszczony na jednym z segmentów, oznaczając miejsce pobytu oddziału.
Poruszanie się po terenie labiryntu następuje zawsze z jednego segmentu do segment sąsiedniego, to znaczy: niedopuszczalne jest przejście od razu dwóch lub więcej segmentów.

\subsection{Sposób postępowania}
Gdy oddział zdecyduje się opuścić segment, jeden z graczy losowo wybiera z pojemnika następny segment i umieszcza go tak, by przylegał do segmentu opuszczanego.
Należy przy tym pamiętać, by korytarze lub przejścia na obu segmentach łączyły się ze sobą.
Jeżeli istnieje kilka sposobów ułożenia nowego segmentu i każdy z nich spełnia sformułowany w poprzednim zdaniu warunek, gracze mają całkowitą swobodę w wyborze jednego z tych sposobów.
Jeżeli oddział porusza się drogą, którą już raz przeszedł, nie istnieje konieczność ciągnięcia nowych segmentów.
